\documentclass[11pt,a4paper]{article}
\usepackage{gtu-el-yazisi}
\begin{document}
\notbilgi{Alan Teorisi}{Ders 15: Lorentz kuvveti, magnetik indüksiyon, akım yoğunluğu}{Ders15.pdf}

\sayfa{1}
\elcizimi[0.3]{sekiller/Ders15-s1-1.png}[1]

\elcizimi[0.55]{sekiller/Ders15-s1-2.png}[1]

\noindent $q_i$: \quad $i = 1,2,\dots,n$

$\vec{E}$ $q_i$'lerin yarattığı alansa, \okunamadi{yüklerde} $q$'ya göre hareket halinde ise o zaman gözlenen $\vec{F}$ kuvveti $q\vec{E}$'den farklı olur ve bu fark $q$'ya ve $q$'nun $\vec{v}$'sine dik olacak şekilde ortaya çıkan magnetik indüksiyon vektörü adı verilir ve bu hareketli yükler tarafından yaratılıyor.

\[ \vec{F} = q\cdot\vec{E} \]

\elcizimi[0.5]{sekiller/Ders15-s1-3.png}[1]

\[ \vec{F} - q\vec{E} = q\,\vec{v}\times\vec{B} \]
\formulkutu[Lorentz kuvveti]{\vec{F} = q\vec{E} + q\,\vec{v}\times\vec{B}}
\[ = q\,(\vec{E} + \vec{v}\times\vec{B}) \qquad\qquad \vec{B}\ ? \]

\noindent $\vec{E}$'yi yaratan kaynaklar durgun yükler

\noindent $\vec{B}$'yi \quad '' \qquad '' \qquad\qquad hareketli yükler

\elcizimi[0.7]{sekiller/Ders15-s1-4.png}[1]

\noindent ``iki ayrı gözlemci''
\begin{align*}
  \vec{v} &= \vec{v}_1 \\
  \vec{F} &= q\,(\vec{E} + \vec{v}\times\vec{B}) \\
  \vec{F}_1' &= q_1\,(\vec{E}' + \vec{v}'\times\vec{B}')
\end{align*}
\textit{(Son satırda $\vec{v}'\times\vec{B}'$ terimi çizilmiş.)}

\[ \begin{aligned} q_i:\quad \vec{v}_i = 0 \qquad & \vec{B} = 0 \\ \vec{v}_i \neq 0 \qquad & \vec{B} \neq 0 \quad (t\text{'den bağımsız}) \end{aligned} \]

\[ \vec{F} = q\,(\vec{E} + \vec{v}\times\vec{B}) = q\vec{E} + q\,\vec{v}\times\underbrace{\vec{B}}_{\substack{\text{Tesla}\\ \text{Weber/m}^2}} \]

\[ \vec{F} = q\,\vec{v}\times\vec{B} \qquad\qquad \vec{B} = 1\,\mathrm{T} \]
\elcizimi[0.3]{sekiller/Ders15-s1-5.png}[1]
\[ |\vec{F}| = |q|\,|\vec{v}|\,|\vec{B}|\cdot\sin\theta \]

\[ \vec{v} \parallel \vec{B} \]
\elcizimi[0.9]{sekiller/Ders15-s1-6.png}[1]
\[ \vec{F} = 0 \]

\begin{align*}
  \vec{F} &= q\,(\vec{v}\times\vec{B} + \vec{E}) \\
  -\vec{F} &= q\,(-\vec{v}\times\vec{B}' - \vec{E}')
\end{align*}
\[ \Downarrow \quad \vec{B} \]

\elcizimi[0.4]{sekiller/Ders15-s1-7.png}[1]

\begin{align*}
  \vec{B}' &\neq -\vec{B} \\
  \vec{E}' &= -\vec{E}\ \checkmark \\
  \vec{v}' &= -\vec{v}\ \checkmark \\
  \vec{F}' &= -\vec{F}\ \checkmark
\end{align*}

\sayfa{2}
\noindent Pr 2 \qquad\qquad $\vec{a}(x,y,z)$ \qquad\qquad $\vec{a}(x',y',z')$

\elcizimi[0.5]{sekiller/Ders15-s2-1.png}[2]

\begin{align*}
  \vec{a}'(x',y',z') &= -\vec{a} \\
  \vec{B} = \operatorname{rot}\vec{a} &= \left(\frac{\partial}{\partial x}\vec{e}_x + \frac{\partial}{\partial y}\vec{e}_y + \frac{\partial}{\partial z}\vec{e}_z\right)\times\vec{a} \\
  \vec{B}' = \operatorname{rot}\vec{a}' &= \left(\frac{\partial}{\partial x'}\vec{e}_{x'} + \frac{\partial}{\partial y'}\vec{e}_{y'} + \frac{\partial}{\partial z'}\vec{e}_{z'}\right)\times\vec{a}' \\
  &= \left(\textcolor{kirmizikalem}{-}\frac{\partial}{\partial x}\vec{e}_x \textcolor{kirmizikalem}{-} \frac{\partial}{\partial y}\vec{e}_y \textcolor{kirmizikalem}{-} \frac{\partial}{\partial z}\vec{e}_z\right)\times\left(\textcolor{kirmizikalem}{-}\vec{a}\right) \\
  &= \operatorname{rot}\vec{a} = \vec{B}
\end{align*}

\noindent $a$ kutupsal ise rot $a$ eksenel vektördür.

\elcizimi[0.22]{sekiller/Ders15-s2-2.png}[2]

\begin{align*}
  dq &= \rho\, dv \\
  d\vec{B} &= dq\,\frac{\vec{v}\times\vec{r}'}{r'^3}
\end{align*}

\elcizimi[0.5]{sekiller/Ders15-s2-3.png}[2]

\[ \vec{B}\perp\vec{r}' \qquad \vec{B}\perp\vec{v} \qquad |\vec{B}| \sim \frac{1}{r'^2}\,q\,|\vec{v}| \]

\[ \vec{B} = q\cdot\frac{\vec{v}\times\vec{r}'}{r'^3} \qquad\qquad \vec{a}\times\vec{b} = \vec{c}\ \text{ise}\ \ \vec{c}\perp\vec{a},\ \vec{c}\perp\vec{b} \]
\[ |\vec{B}| = |q|\,|\vec{v}|\cdot\frac{1}{r'^2}\cdot\sin\theta \]

\noindent yük yoğunluğu $\rho$ ile dağılmışsa

\begin{align*}
  dq &= \rho\, dV \\
  d\vec{B} &= dq\,\frac{\vec{v}\times\vec{r}'}{r'^3} \\
  &= \rho\,\vec{v}\times\frac{\vec{r}'}{r'^3}\, dV \\
  \vec{B} &= \int_V d\vec{B} = \int_V \rho\,\vec{v}\times\frac{\vec{r}'}{r'^3}\, dV
\end{align*}

\elcizimi[0.45]{sekiller/Ders15-s2-4.png}[2]

\[ \underbrace{\rho}_{\mathrm{C/m^3}}\cdot\underbrace{\vec{v}}_{\mathrm{m/s}} \equiv \underbrace{\vec{J}}_{\text{hacimsel akım yoğunluğu}}\ (\mathrm{A/m^2}) \qquad\qquad \mathrm{A} = \mathrm{C/s} \]

\[ \vec{B} = \frac{\mu_0}{4\pi}\int \vec{J}\times\frac{\vec{r}'}{r'^3}\, dv \qquad \vec{J}\parallel\vec{v} \qquad \vec{B}\perp\vec{r}',\ \vec{B}\perp\vec{J} \]
\[ \vec{E} = \frac{1}{4\pi\varepsilon_0}\int \frac{\rho\,\vec{r}'}{r'^3}\, dv \]

\sayfa{3}
\elcizimi[0.4]{sekiller/Ders15-s3-1.png}[3]

\elcizimi[0.55]{sekiller/Ders15-s3-2.png}[3]

\[ t_0,\ t_0 + \Delta t \]
\[ \lim_{\Delta t\to 0}\frac{\Delta q}{\Delta t} \equiv I_S(t) \qquad \text{($S$'den geçen elektrik akımının tanımı)} \]

\elcizimi[0.85]{sekiller/Ders15-s3-3.png}[3]

\[ \vec{J} = \rho\vec{v} \]

\[ I_{\Delta S} = \frac{\vec{v}\cdot\vec{n}\,\Delta t\,\Delta S\,\rho}{\Delta t} = \vec{J}\cdot\vec{n}\,\Delta S \]

\begin{align*}
  dq &= \rho\, dV \\
  \Delta q &= \rho\cdot\underbrace{\vec{v}\,\Delta t\cdot\vec{n}\,\Delta S}_{dV}
\end{align*}

\elcizimi[0.45]{sekiller/Ders15-s3-4.png}[3]

\elcizimi[0.5]{sekiller/Ders15-s3-5.png}[3]

\[ \vec{n}\cdot\underbrace{\vec{h}'}_{\vec{v}\cdot\Delta t} = |\vec{n}|\,|\vec{h}'|\cdot\cos\theta = h'\cos\theta \]

\[ \lim_{t\to 0}\frac{\Delta q}{\Delta t} \equiv I_S(t) \]

\[ dI = \vec{J}\cdot d\vec{S} \]
\elcizimi[0.35]{sekiller/Ders15-s3-6.png}[3]
\[ \checkmark\ \ I = \int_S dI = \int_S \vec{J}\cdot d\vec{S} = \int_S \vec{J}\cdot\vec{n}\, dS \qquad d\vec{S} = dS\cdot\vec{n} \]

\noindent akım skalerdir

\elcizimi[0.4]{sekiller/Ders15-s3-7.png}[3]

\elcizimi[0.4]{sekiller/Ders15-s3-8.png}[3]

\[ J\ [\mathrm{A/m^2}] \]

\noindent hacim içinden geçen\\
Akım diye bir şey yoktur\\
yüzey içinden geçen akım vardır

\section{3.2.1}
\[ H = \frac{1}{4\pi}\int_{-\infty}^{\infty}\left(\mathbf{e}_z\times\frac{\vec{r}'}{r'}\right)\frac{d\xi}{r'^2} = \frac{1}{2\pi\rho}\,\mathbf{e}_\varphi \]
\[ \vec{B} = \mu_0\vec{H} \quad \text{(boşlukta)} \]

\sayfa{4}
\[ [\vec{H}] = \mathrm{A/m} \qquad \mu_0 = 4\pi\cdot 10^{-7}\ \mathrm{H/m} \]
\[ [\vec{B}] = \underbrace{\mathrm{Wb}}_{\text{A.Henry}}/\mathrm{m}^2 \equiv \mathrm{T} \]

\elcizimi[0.8]{sekiller/Ders15-s4-1.png}[4]

\[ r' = |\vec{r}'| = \sqrt{(x-\xi)^2 + (y-\eta)^2 + (z-\zeta)^2} \]
\[ \vec{r}' = (x- \]

\begin{align*}
  \vec{J} &= I\,\delta(\xi)\,\delta(\eta)\,\vec{e}_z \\
  \vec{B} &= \mu_0\vec{H} \;\Rightarrow\; \vec{H} = \frac{1}{\mu_0}\vec{B}
\end{align*}

\[ = \frac{1}{4\pi}\iiint I\,\delta(\xi)\,\delta(\eta)\,\vec{e}_z\times\frac{(x-\xi)\vec{e}_x + (y-\eta)\vec{e}_y + (z-\zeta)\vec{e}_z}{\left[(x-\xi)^2 + (y-\eta)^2 + (z-\zeta)^2\right]^{3/2}}\, d\xi\, d\eta\, d\zeta \]

\[ \int_{-\infty}^{\infty}\delta(\xi)\,f(\xi)\,d\xi = f(0) \]
\[ \Rightarrow\ \vec{H} = \frac{I}{4\pi}\int_{\zeta=-\infty}^{\infty}\frac{x\,\vec{e}_y - y\,\vec{e}_x}{\left[x^2+y^2+(z-\zeta)^2\right]^{3/2}}\, d\zeta \]

\begin{align*}
  \vec{H} &= \frac{I}{4\pi}\,(x\vec{e}_y - y\vec{e}_x)\int_{-\infty}^{+\infty}\frac{d\zeta}{\left(\sqrt{\rho^2+(z-\zeta)^2}\right)^3} \\
  \vec{H} &= \frac{I\,(x\vec{e}_y - y\vec{e}_x)}{4\pi}\int_{\psi=+\pi/2}^{-\pi/2}\frac{-\rho\,(1+\tan^2\psi)\,d\psi}{(\rho^2+\rho^2\tan^2\psi)^{3/2}} \\
  &= \frac{I\rho\,(x\vec{e}_y - y\vec{e}_x)}{4\pi\rho^3}\int_{-\pi/2}^{\pi/2}\frac{d\psi}{\sqrt{1+\tan^2\psi}} \\
  &= \frac{I}{4\pi\rho^2}\,\bigl(\underbrace{\rho\cos\phi\,\vec{e}_y - \rho\sin\phi\,\vec{e}_x}_{\vec{e}_\phi}\bigr)\int_{-\pi/2}^{\pi/2}\cos\psi\, d\psi
\end{align*}

\[ \sqrt{1+\tan^2\psi} = \sqrt{\frac{1}{\cos^2\psi}} = \frac{1}{\cos\psi} \]

\begin{align*}
  &= \frac{I}{4\pi\rho}\,\vec{e}_\phi\cdot\sin\psi\Big|_{-\pi/2}^{\pi/2} \\
  &= \frac{I}{2\pi\rho}\,\vec{e}_\phi
\end{align*}

\[ \rho = \sqrt{x^2+y^2} \]

\elcizimi[0.35]{sekiller/Ders15-s4-2.png}[4]

\begin{align*}
  \tan\psi &= \frac{z-\zeta}{\rho} \\
  \rho\,(1+\tan^2\psi)\,d\psi &= -d\zeta \\
  d\zeta &= -\rho\,(1+\tan^2\psi)\,d\psi
\end{align*}

\sayfa{5}
\elcizimi[0.7]{sekiller/Ders15-s5-1.png}[5]

\elcizimi[0.4]{sekiller/Ders15-s5-2.png}[5]

\[ \vec{e}_\phi = \cos\phi\,\vec{e}_y - \sin\phi\,\vec{e}_x \]

\[ \begin{aligned} V &\longleftrightarrow \vec{A} \\ \rho &\longleftrightarrow \vec{J} \\ \varepsilon_0 &\longleftrightarrow \frac{1}{\mu_0} \end{aligned} \qquad\qquad \begin{aligned} \Delta V &= -\rho/\varepsilon_0 \\ \Delta\vec{A} &= -\mu_0\vec{J} \end{aligned} \]
\[ \left\{ \begin{aligned} \operatorname{rot}\vec{E} &= 0 \\ \operatorname{div}\vec{B} &= 0 \end{aligned} \right. \]

\[ \Delta A = -\mu_0 J \]
\[ \kutula{\vec{H} = \frac{1}{\mu_0}\operatorname{rot}\vec{A}} \]

\noindent $\operatorname{div}A = 0$ \qquad rot \okunamadi{}

\noindent Pr 3

\elcizimi[0.5]{sekiller/Ders15-s5-3.png}[5]

\elcizimi[0.4]{sekiller/Ders15-s5-4.png}[5]

\begin{align*}
  \vec{B} &= \frac{\mu_0}{2\pi}\,\frac{I}{\rho}\,\vec{e}_\phi \\
  \vec{B}_q &= \frac{\mu_0 I}{2\pi a}\,\vec{e}_y \\
  \vec{F}_q &= q\vec{v}_0\times\vec{B}_q \\
  &= q v_0\,\vec{e}_z\times\frac{\mu_0 I}{2\pi a}\,\vec{e}_y \\
  \vec{F}_q &= -\frac{\mu_0\, q\, v_0\, I}{2\pi a}\,\vec{e}_x \\
  \vec{F}_I &= -\frac{\mu_0\, q\, v_0\, I}{2\pi a}\,\vec{e}_x \\
  \vec{F}_{I1} &= +\frac{\mu_0\, q\, v_0\, I_1}{2\pi\,(b-a)}\,\vec{e}_x
\end{align*}

\begin{align*}
  \vec{F}_q &= \vec{F}_I + \vec{F}_{I1} \\
  &= \frac{\mu_0\, q\, v_0}{2\pi}\,\vec{e}_x\left(\frac{I_1}{b-a} - \frac{I}{a}\right)
\end{align*}

\[ \vec{F}_q = 0 \;\Longrightarrow\; \frac{I_1}{b-a} - \frac{I}{a} = 0 \]
\[ \kutula{I_1 = \frac{I\,(b-a)}{a}} \]

\sayfa{6}
\noindent Pr 10

\[ \vec{F} = \underbrace{I_2\,\vec{l}}_{q\vec{v}} \times \underbrace{\vec{B}_T}_{\frac{\mu_0}{2\pi}\,\vec{e}_x\left(\frac{I_3}{b}-\frac{I_1}{a}\right)} \]

\begin{align*}
  \vec{B}_1 &= \frac{\mu_0 I_1}{2\pi a}\,(-\vec{e}_x) \\
  \vec{B}_3 &= \frac{\mu_0 I_3}{2\pi b}\,\vec{e}_x
\end{align*}

\begin{align*}
  \vec{F} &= (I_2\, l\,\vec{e}_z)\times\left(\frac{\mu_0}{2\pi}\,\vec{e}_x\left(\frac{I_3}{b}-\frac{I_1}{a}\right)\right) \\
  &= \frac{I_2\, l\,\mu_0}{2\pi}\left(\frac{I_3}{b}-\frac{I_1}{a}\right)\vec{e}_y
\end{align*}

\[ \vec{F} = 0 \ \text{için} \qquad \frac{I_3}{b} = \frac{I_1}{a} \]

\[ \begin{aligned} &q\vec{v}\,dt\times\vec{B} \\ &\vec{J}\times\vec{B} \\ &\ \| \\ &I\cdot\vec{l} \end{aligned} \qquad\qquad \vec{F} = \int_V(\vec{J}\times\vec{B})\,dv \]

\noindent Sene başından geldiğimiz yere kadar \qquad 9:30 \qquad 3 soru \qquad \okunamadi{Ağu 25}:

\noindent Birtane bundan\\
Pr dikkate edin\\
Kalan problemler\\
3. ödevi \okunamadi{Aslıhan}

\end{document}
